\section{Schedule rigidity and upper-envelope flexibility}
\label{sec:schedules}

The asymptotic law does not license arbitrary exact convergence schedules.  The integer geometry of the positions creates a strong restriction, while the freedom to enlarge future gaps gives a complementary one-sided realization result.

\begin{theorem}[Rigidity/flexibility dichotomy]
\label{thm:schedules}
The following statements hold.
\begin{enumerate}
    \item For every increasing-gap sequence,
    \begin{equation}
    n_{j+1}\geq 1+jq_1+\frac{j(j-1)}{2}.
    \label{eq:quadratic-position}
    \end{equation}
    There is a constant $C_\beta>0$ such that
    \begin{equation}
    0<L_j\leq C_\beta\beta^{-n_{j+1}}.
    \label{eq:superexp-bound}
    \end{equation}
    In particular, $j^rL_j\to0$ for every $r>0$.
    \item Let $(\varepsilon_j)_{j\geq2}$ be any sequence of positive real numbers.  There exists a strictly increasing integer gap sequence for which
    \begin{equation}
    0<L_j\leq\varepsilon_j\qquad(j\geq2).
    \label{eq:upper-envelope}
    \end{equation}
\end{enumerate}
The second assertion is an upper-envelope statement; it does not assert $L_j\sim\varepsilon_j$ or equality at any index.
\end{theorem}

\begin{proof}
Because the $q_j$ are strictly increasing positive integers,
$q_j\geq q_1+j-1$.  Summing \eqref{eq:positions} gives \eqref{eq:quadratic-position}.  Moreover, $M_j\geq\beta^{-1}$ and \eqref{eq:tail-geometric} imply
\[
 L_j\leq\frac{T_j}{M_j}
 \leq\frac{\beta\,\beta^{-n_{j+1}}}{1-\beta^{-q_{j+1}}}
 \leq\frac{\beta}{1-\beta^{-1}}\,\beta^{-n_{j+1}}.
\]
Thus \eqref{eq:superexp-bound} holds with $C_\beta=\beta/(1-\beta^{-1})$.  Combining it with \eqref{eq:quadratic-position} proves the final assertion of part~1.

For part~2, fix $q_1=2$, so $n_1=1$, $n_2=3$, and let $\underline\beta>1$ be the solution of
\[
 \underline\beta^{-1}+\underline\beta^{-3}=1.
\]
Any infinite completion has $\beta\geq\underline\beta$.  Suppose $q_1,\ldots,q_{j-1}$ have been chosen.  Choose an integer $q_j>q_{j-1}$ so large that, with $n_{j+1}=n_j+q_j$,
\begin{equation}
 \frac{2\underline\beta^{-n_{j+1}}}{1-\underline\beta^{-1}}
 \leq\varepsilon_j.
 \label{eq:recursive-choice}
\end{equation}
This recursion defines an infinite strictly increasing sequence.

After completing it, use $\beta<2$ to obtain $M_j\geq\beta^{-1}>1/2$.  Since future positions are distinct integers at least $n_{j+1}$ and $\beta\geq\underline\beta$,
\[
 T_j\leq\sum_{n\geq n_{j+1}}\underline\beta^{-n}
 =\frac{\underline\beta^{-n_{j+1}}}{1-\underline\beta^{-1}}.
\]
The upper half of \eqref{eq:certificate}, followed by \eqref{eq:recursive-choice}, gives $L_j\leq\varepsilon_j$.
\end{proof}

\begin{remark}[Why exact schedules are excluded]
Taking $\varepsilon_j=1/j$ illustrates the obstruction: \cref{thm:schedules}(1) forces $L_j=o(j^{-r})$ for every $r$, so $L_j\sim1/j$ is impossible.  Even among faster schedules, \cref{cor:log-rigidity} confines the principal scale to the discrete locations $\beta^{-n_{j+1}}$.  Characterizing which two-sided target sequences are realizable would require controlling the feedback of the entire chosen tail on $\beta$ and $D_\beta$; that problem is not solved here.
\end{remark}
